\documentclass[10pt,a4paper]{amsart}
\usepackage[margin=19mm]{geometry}
\usepackage{amsmath,amssymb,mathtools}
\usepackage[scr=rsfs,cal=euler]{mathalfa}
\usepackage{xfrac}
\usepackage[unicode,hidelinks,bookmarks=false]{hyperref}
\newcommand{\ie}{i.e.\ }
\newcommand{\cf}{cf.\ }
\newcommand{\resp}{respectively\ }
\newcommand{\Subsection}{\subsection}
\providecommand{\nobreakdash}{\nobreak\hbox{-}}
\setlength{\emergencystretch}{2em}
\multlinegap0pt
\usepackage[shortlabels]{enumitem}
\usepackage{xcolor}
\usepackage{amssymb}
\usepackage{dsfont}
\usepackage{mathtools}
\usepackage[normalem]{ulem}
\usepackage{csquotes}
\newtheorem{assumption}{Assumption}
\newtheorem{proposition}{Proposition}
\newtheorem{theorem}{Theorem}
\usepackage{cleveref}
\crefname{assumption}{Assumption}{Assumptions}
\crefname{proposition}{Proposition}{Propositions}
\crefname{theorem}{Theorem}{Theorems}
\newcommand{\N}{\mathbbm{N}}
\newcommand{\R}{\mathbbm{R}}
\newcommand{\fp}{\mathfrak{p}}
\newcommand{\fq}{\mathfrak{q}}
\providecommand{\mathbbm}{\mathbb} % In case we don't load bbm
\newcommand{\placeholder}{\mathord{\color{black!33}\bullet}}
\newcommand{\sD}{\mathsf{D}}
\title{Operator Learning for Smoothing and Forecasting}
\author{Edoardo Calvello, Elizabeth Carlson, Nikola Kovachki, Michael N. Manta, Andrew M. Stuart}
\date{Source: arXiv:2603.20359v2 (CC BY 4.0)}
\begin{document}
\maketitle

\noindent\emph{Source and licence.} Operator Learning for Smoothing and Forecasting. Version arXiv:2603.20359v2 (CC BY 4.0), distributed by arXiv under the Creative Commons Attribution 4.0 International licence, http://creativecommons.org/licenses/by/4.0/, as recorded in the arXiv licence metadata for this exact version and shown on the abstract page https://arxiv.org/abs/2603.20359v2. Full licence notice: \texttt{licenses/arxiv-2603.20359-CC-BY-4.0.txt}.\par\smallskip

\noindent\textit{Editorial scope.} Counted unit: the article's theorem thm:UAforecasting (source0.tex lines 916--921). The article's complete original proof of it is delivered with it (source0.tex lines 922--949, one \texttt{proof} environment of 28 lines). No mathematics is written, repaired or re-derived: the statement and the proof are the article's own lines, in the article's own notation, and no retained line is substituted or re-flowed.  Retained uncounted as the source-local context the proof needs: lines 386--402; lines 571--579; lines 729--732; lines 810--820; lines 895--905.  Nothing was omitted from any retained span, so the package carries no omission record.  No retained line contains a citation, so the wrapper carries no bibliography and no bibliography entry.  No retained line contains a figure environment, an image-inclusion command or drawing code, so no image is delivered and the package declares no asset dependency.  Editorial formatting only, in addition: an AMS \texttt{amsart} preamble that restates the article's own declarations and macros used by the retained lines, namely the environments \texttt{assumption}, \texttt{proposition}, \texttt{theorem} on separate counters, together with the standard packages those lines need.  The retained environments print in this document as Assumption 1, Proposition 1, Theorem 1 in order of appearance, whereas the article prints the same items with the numbering of its own full document.  The complete native source of the article as distributed in the arXiv e-print for this exact version is bundled unchanged as \texttt{sources/196840/source0.tex}.
Machine learning has opened new frontiers in data assimilation and forecasting in dynamical systems,
ranging from methods which improve on traditional model-driven 
algorithms, through to new purely data-driven methods. The goal
of this paper is to establish a mathematical approach to the study
of the purely data-driven methods. To this end, we consider the general dynamical system given by
\begin{subequations}
\label{eq:dynamical_system_intro}
\begin{align}
        \label{eq:dynsys_f}
        \dot{p} &= f(p, q), \qquad p(0) = p_0 \in \R^{d_p}, \\
        \label{eq:dynsys_g}
        \dot{q} &= g(p, q), \qquad q(0) = q_0 \in \R^{d_q}.
\end{align}
\end{subequations}
Our focus is on two problems: (i) existence, universal approximation, and approximation from data
of the map from observed $\{p(t)\}_{t\in [0,T]}$ to unobserved $\{q(t)\}_{t\in [0,T]}$; (ii)
existence, universal approximation, and approximation from data of the map from  observed $\{p(t)\}_{t\in [0,T]}$ to unobserved $\{p(t)\}_{t\in [T,T+\tau]}$.

\begin{assumption}[Regularity]
\label{assump:regularity}
 Assume that, for some $k \in \N$,
 \begin{itemize}
     \item $f \in C^k (\R^{d_p + d_q}; \R^{d_p})$;
     \item $g \in C^k (\R^{d_p + d_q}; \R^{d_q})$.
 \end{itemize}
Assume, furthermore, that solutions to \eqref{eq:dynamical_system_intro} exist for all $t \in \R^+.$
\end{assumption}

\begin{assumption}[Observability-Rank Condition]
\label{assump:observability}
    We say that the system \eqref{eq:dynamical_system_intro} satisfies the observability-rank condition at $(\mathfrak{p}, \mathfrak{q})\in \R^{d_p + d_q}$, if there exists $n\in\N$ and a linear operator $L\colon\R^{(n+1)d_p}\to\R^{d_p+d_q}$ such that the rank of the matrix $\sD LF^{(n)}\bigl(\fp, \fq \bigr)$ is $d_p + d_q$.
\end{assumption}

\begin{proposition}
\label{thm:UANeuralOperator}
Let $D \subset \R^d$ be a bounded domain with Lipschitz boundary, and fix integers $s,s'\geq 0$. If $
\Psi^\dagger\colon C^{s}(\overline{D};\,\R^{r}) \to C^{s'}(\overline{D};\,\R^{r'})$
is a continuous operator and $K \subset C^{s}(\overline{D};\,\R^{r})$ a compact set, then for any $\varepsilon>0$ there exists a transformer neural operator 
$\Psi(\placeholder; \theta)\colon K \subset C^{s}(\overline{D};\,\R^{r}) \to C^{s'}(\overline{D};\,\R^{r'})$
so that
\begin{equation}
\sup_{u\in K} \left\| \Psi^\dagger(u) - \Psi(u;\theta) \right\|_{C^{s'}} \le \varepsilon .
\end{equation}
\end{proposition}

\begin{proposition}[Existence of map $\Psi_F^\dagger:p \!\restriction_{[0,T]}\mapsto p \!\restriction_{[T,T+\tau]}$]
\label{prop:psiF}
Consider the dynamical system \eqref{eq:dynamical_system_intro} satisfying 
regularity \Cref{assump:regularity}, for integer $k$, and the observability \Cref{assump:observability} at some $(\mathfrak{p},\mathfrak{q})\in\R^{d_p + d_q}$, for integer $n \le k.$
Let $ U \ni (\fp, \fq)$ and $V \ni LF^{(n)}(\fp, \fq)$ be the open sets given by the inverse function theorem. 
Then, for every compact $I \subset U$, there exists a continuous operator $\Psi_F^\dagger \colon S^I_{[0,T],p} \subset C^k([0,T];\R^{d_p}) \to S^I_{[T,T+\tau],p} \subset C^k([T, T+\tau]; \R^{d_p})$ such that, for every $(p, q) \in S^I_{[0,T]}$,
\begin{equation}
    \label{eq:p_to_q_map3}
    \Psi_F^\dagger(p)(t) = p(t), \qquad T\leq t\leq T+\tau.
\end{equation}
\end{proposition}

\begin{theorem}[Universal Approximation for Forecasting]
\label{thm:UAforecasting}
Consider the dynamical system \eqref{eq:dynamical_system_intro} satisfying
regularity \Cref{assump:regularity}, for integer $k$, and the observability \Cref{assump:observability} at some $(\mathfrak{p},\mathfrak{q})\in\R^{d_p + d_q}$, for integer $n \le k.$ For any $\epsilon > 0$ there exists a transformer neural operator $\Psi(\placeholder;\theta)\colon S^I_{[0,T],p} \subset C^k ([0,T];\R^{d_p}) \to S^I_{[T,T+\tau],p} \subset C^k ([T,T+\tau];\R^{d_p})$ satisfying
\[\sup_{p \in S^I_{[0,T],p}} \|\Psi_F^\dagger(p) - \Psi(p;\theta)\|_{C^k([T,T+\tau];\R^{d_p})} \leq \epsilon. \]
\end{theorem}

\begin{proof}
    We know that $\Psi_F^{\dagger}$ exists and is continuous since the assumptions of \Cref{prop:psiF} are satisfied. Recall that $S^I_{[0,T],p}$ is compact. We apply the result of \Cref{thm:UANeuralOperator} to an operator defined via linear homeomorphisms of $\Psi_F$, which we will then use to establish the approximation result in its final form. Indeed, we define the translation map $\Lambda_T:C^k([T,T+\tau];\R^{d_p})\to C^k([0,\tau];\R^{d_p})$ so that
    \[(\Lambda_T \varphi)(s) = \varphi(T+s), \qquad s\in[0,\tau], \]
    for any $\varphi\in C^k([T,T+\tau];\R^{d_p})$. We note that 
    \[\|\varphi\|_{C^k([T,T+\tau];\R^{d_p})} = \|\Lambda_T\varphi\|_{C^k([0,\tau];\R^{d_p})}.\]
    It is also possible to define the rescaling operator $R_\tau:C^k([0,\tau];\R^{d_p})\to C^k([0,1];\R^{d_p})$ so that
    \[
    (R_\tau\varphi)(s) = \varphi(\tau s), \qquad s\in[0,1],
    \]
    for any $\varphi\in C^k([0,\tau];\R^{d_p})$. We similarly define the rescaling operator $R_T$. We note that these rescaling operators are linear homeomorphisms, preserving equivalence of norms up to multiplicative constants. We therefore define the translated and rescaled forecasting operator $\widetilde{\Psi}^\dagger_F:R_T(S^I_{[0,T],p})\subset C^k([0,1];\R^{d_p})\to C^k([0,1];\R^{d_p})$ defined via the composition
    \[
    \widetilde{\Psi}^\dagger_F = R_\tau \circ \Lambda_T \circ {\Psi}^\dagger_F \circ R_T^{-1},
    \]
    which is continuous by continuity of all the operators in the composition. We note that $R_T(S^I_{[0,T],p})$ is compact since $R_T$ is a homeomorphism and therefore preserves compactness. Hence, by \Cref{thm:UANeuralOperator}, it is possible to deduce that for any $\delta>0$ there exists a neural operator $\widetilde{\Psi}(\placeholder,\theta): R_T(S^I_{[0,T],p})\subset C^k([0,1];\R^{d_p})\to C^k([0,1];\R^{d_p})$ satisfying
    \begin{equation}
    \label{eq:approximant}
    \sup_{p \in R_T(S^I_{[0,T],p})} \|\widetilde{\Psi}_F^\dagger(p) - \widetilde{\Psi}(p;\theta)\|_{C^k([0,1];\R^{d_p})} \leq \delta. 
    \end{equation}
    This indeed follows from a direct application of \Cref{thm:UANeuralOperator} with $D = [0, 1]$, $r = d_p$, $r' = d_p$, $\Psi^{\dagger} = \widetilde{\Psi}_F^{\dagger}$, and $K = R_T(S^I_{[0,T],p})$. Defining $\Psi(\placeholder,\theta):C^k ([0,T];\R^{d_p}) \to C^k ([T,T+\tau];\R^{d_p})$ via the composition
    \[
    \Psi(\placeholder,\theta) = \Lambda_T^{-1}\circ R_\tau^{-1} \circ \widetilde{\Psi}(\placeholder;\theta)\circ R_T,
    \]
    it is readily deduced that for any $p \in S^I_{[0,T],p}$ it holds that
    \[
    \|\Psi_F^\dagger(p) - \Psi(p;\theta)\|_{C^k([T,T+\tau];\R^{d_p})} = C\|\widetilde{\Psi}_F^\dagger(R_Tp) - \widetilde{\Psi}(R_Tp;\theta)\|_{C^k([0,1];\R^{d_p})},
    \]
    for a constant $C$ depending on $T$ and $\tau$. Therefore, choosing $\epsilon=\delta/C$ and applying the result established in \eqref{eq:approximant} yields the conclusion.
    \end{proof}

\end{document}
